% =====================================================================
%  23.1  The Li Chao Line Container.
%  Exports: sec:c29-lichao  def:c29-lichao  thm:c29-lichao-cost
%           fig:c29-lichao  lst:c29-lichao  prob:c29-631e  lst:c29-631e
%  Solved CF 631E "Product Sum" (2600), live-verified via the problemset
%  API (tags data structures / dp / geometry).  Official samples read live:
%    S1: n=4, a=(4,3,2,5) -> 39 ; S2: n=5, a=(1,1,2,7,1) -> 49 ; S3: n=3,
%    a=(1,1,2) -> 9.  Hand-verified here: S1 original characteristic
%    S = 1*4+2*3+3*2+4*5 = 36; best single move gain = 3 (move a_1=4 from
%    index 1 to index 3, or a_3=2 from index 3 to index 1) -> 39.
% =====================================================================
\section{The Li Chao Line Container}\label{sec:c29-lichao}

\begin{previouslybox}[Ch.~19 indexed positions and folded a fixed array; what if the ``array'' is a moving family of lines?]
A Chapter~19 segment tree is built over \emph{positions} \(1..n\) and folds a
value stored at each.  Now ask something the position-tree cannot: keep a set of
lines \(\ell(x)=mx+b\), and for any query point \(x\) return the one that is
smallest (or largest) there.  The lines are inserted \emph{online}, and which
line wins depends on \(x\) -- there is no fixed array to fold.  If the slopes
arrive in monotone order a convex-hull deque suffices, but slopes of arbitrary
sign and size kill that.  The \term{Li Chao line container} answers it anyway,
in \(\bigO{\log C}\) per insert and per query, where \(C\) is the size of the
(integer) domain of \(x\).
\end{previouslybox}

\begin{topicintro}
A Li Chao container is a segment tree whose axis is the \emph{\(x\)-domain}, not
the array of positions.  Each node owns an interval of \(x\)-values and stores
\emph{one line}.  To insert a line, compare it with the stored line at the
node's midpoint: keep the better one at the node, and recurse the loser to the
\emph{single} side on which it can still beat the winner.  Because two lines
cross at most once, only one of the two halves can contain that crossing, so
each insert descends one root-to-leaf path -- \(\bigO{\log C}\).  A point query
walks that leaf's path and takes the best of the stored lines.  This is the tool
for a DP transition \(f(i)=\min_{j}\{m_j\,x_i+b_j\}+c_i\) whose slopes \(m_j\)
are \emph{arbitrary}: each earlier state contributes a line, each new state is a
query.

\textbf{Prerequisites.} Chapter~19 (\chref{c13}) -- the segment-tree skeleton and
the idea of an interval-owning node.  A line evaluated at a point, and the fact
that two distinct lines meet at most once, are the only geometry needed.

\textbf{After this section you can}
\begin{itemize}[nosep]
  \item insert a line into a container over \([1,C]\) by the mid-winner descent
        and read a point query as a root-to-leaf best; \hfill\lv{3}
  \item state and prove the \(\bigO{\log C}\) insert/query bound from ``two lines
        cross once, so the loser can only win on one side''; \hfill\lv{4}
  \item recognise a DP whose transition is an upper/lower envelope of lines with
        arbitrary slopes and wire it to the container; \hfill\lv{4}--\lv{5}
  \item solve CF~631E: sweep destinations, query the best source-line, and get
        the best single-move gain in \(\bigO{n\log C}\). \hfill\lv{5}
\end{itemize}
\end{topicintro}

\begin{recognitionbox}[Pattern card -- ``best of many lines, queried at a point'']
\reckey{a running set of functions each of the form \(m x + b\) (or a value you
can rewrite that way); operations ``add one'' and ``what is the extreme at the
specific \(x\) I hand you''; or a DP whose cost from state \(j\) to state \(i\)
is linear in a quantity \(x_i\) that depends only on the query, with a
\(j\)-dependent intercept -- and \emph{slopes \(m_j\) in no particular order}.}
\recuse{
  \item pointwise \(\min/\max\) over lines with arbitrary slopes -- the container (\(\secref{c29-lichao}\));
  \item that query on a tree, over the ancestor lines -- the rollback container (\(\secref{c29-path}\));
  \item a DP transition \(f(i)=\min_j m_j x_i + b_j + c_i\), \(m_j\) unordered -- insert line \(j\), query at \(x_i\).
}
\recnotuse{
  \item slopes arrive monotone \emph{and} queries monotone -- a convex-hull deque is simpler and \(O(1)\) amortised;
  \item there is a fixed array to fold, no per-\(x\) choice -- Chapter~19's tree (\chref{c13});
  \item the hull must lose lines over time (deletions / arbitrary order online) -- dynamic hull, Ch.~25 (\chref{c23}).
}
One-line test: \emph{insert lines, query the best at an \(x\) -- slopes not
monotone -- index the \(x\)-domain, not positions, and descend once per insert.}
\end{recognitionbox}

% ---------------------------------------------------------------
\subsection{The container and the insertion descent}
% ---------------------------------------------------------------

Fix an integer domain \(x\in[1,C]\).  Build an (implicit) segment tree over that
domain; each node covers an interval \([L,R]\) with midpoint \(M=\lfloor(L+R)/2\rfloor\)
and holds \emph{one} line (initially the \(+\infty\) line for a minimum
container).

\begin{definition}{The Li Chao insert}{c29-lichao}
To insert a new line \(g\) at a node storing line \(f\), covering \([L,R]\) with
mid \(M\): evaluate both at \(L,M,R\).  First ensure \(f\) is the better of the
two at \(M\) (if \(g(M)\) beats \(f(M)\), swap them -- the node keeps the
mid-winner).  Now compare at \(L\): if \(f\) and \(g\) differ in which is better
at \(L\) versus at \(M\), then \(g\) can only beat \(f\) on the \emph{left}
half, so recurse \(g\) into the left child; otherwise if \(g\) beats \(f\) at
\(R\) it can only win on the \emph{right} half, recurse \(g\) into the right
child; otherwise \(g\) is dominated on the whole interval and is dropped.
\end{definition}

\begin{figure}[H]
\centering
\begin{tikzpicture}[diagram, scale=0.85, transform shape]
  \draw[->,darkgray] (0,0) -- (6.6,0) node[below right] {\(x\)};
  \draw[->,darkgray] (0,-1.2) -- (0,3.4) node[above left] {\(y\)};
  \draw[line width=1.3pt, draw=primaryblue]  (0.4,3.0) -- (6.2,-0.4) node[font=\scriptsize,text=primaryblue,anchor=west] at (6.2,-0.4) {\(f\) (kept)};
  \draw[line width=1.3pt, draw=warningred, densely dashed] (0.4,-0.2) -- (6.2,2.9) node[font=\scriptsize,text=warningred,anchor=west] at (4.7,2.3) {\(g\) (loser)};
  \filldraw (2.9,0) circle (1.6pt);
  \node[font=\scriptsize, anchor=north] at (2.9,0) {\(M\): crossing};
  \draw[primaryblue!50, line width=6pt, opacity=0.25] (0.4,-0.2) -- (2.9,0);
  \node[font=\scriptsize, text=warningred, anchor=south east] at (1.5,0.35) {\(g<f\) left only};
  \draw[<-, >=stealth, difforange, line width=1pt] (2.9,-0.9) to[out=-20,in=200] (4.8,-0.9) node[right, font=\scriptsize, text=difforange] {recurse left child};
\end{tikzpicture}
\caption{One insertion step.  The node keeps \(f\), the line that wins at the
midpoint \(M\).  The new line \(g\) loses at \(M\); since two lines cross at
most once, \(g\) can beat \(f\) on \emph{only one side} of \(M\) -- here the
left, where the crossing forces it below.  So \(g\) is recursed to that one
child and the other half is untouched.}
\label{fig:c29-lichao}
\end{figure}

\begin{theorem}{Li Chao insert/query cost}{c29-lichao-cost}
Over a domain of size \(C\), both insertion and point query of a Li Chao
container take \(\bigO{\log C}\) time and the container uses \(\bigO{\log C}\)
new nodes per inserted line (or \(\bigO{C}\) if the tree is pre-built).  A point
query returns the true minimum over all inserted lines.
\end{theorem}
\begin{proofbox}
\textbf{Correctness.} At each visited node the stored line is the better of the
two at that node's midpoint (\(\defref{c29-lichao}\)), so for any \(x\), the line
that is minimal at \(x\) is stored at \emph{some} node on the root-to-leaf path
of \(x\): when a line \(g\) loses the mid at a node it is sent to the unique
half that contains the point where it could beat the kept line, so it is never
discarded before the region where it is genuinely best.  Reading the leaf's whole
path and taking the minimum therefore sees every candidate line that could be
optimal at \(x\).\\
\textbf{Cost.} The tree has depth \(\bigO{\log C}\) over the domain; an insert
follows a single root-to-leaf descent (one recursive call per level, since the
loser goes to exactly one child) and a query likewise follows one path, so both
are \(\bigO{\log C}\). \(\Box\)
\end{proofbox}

\begin{insightbox}[The whole trick is ``two lines cross once'']
Everything hinges on a linearity fact: distinct \(f,g\) satisfy \(f(x)=g(x)\) at
exactly one \(x\), so on either side of that crossing the better line is fixed.
After fixing the winner at the midpoint, the loser can be strictly better on at
most \emph{one} of the two halves -- which is precisely the half straddling the
crossing.  That single observation collapses a would-be ``keep all lines''
search to one descent.  The same fact is why a monotone-slope hull (deque) also
works, but Li Chao needs no ordering of the slopes.
\end{insightbox}

\begin{lstlisting}[caption={Li Chao line container over the integer domain [lo,hi]: insert a line keeping the mid-winner, query the pointwise minimum along the root-to-leaf path. ASCII only}, label={lst:c29-lichao}]
struct Line { long long m, b;                 // y = m*x + b
    long long get(long long x) const { return m*x + b; } };
const long long INF = (1LL<<62);
Line tree_[4*MAXN]; bool has_[4*MAXN];

void insert(int i, int lo, int hi, Line g){   // insert g into node i = [lo,hi]
    int mid = (lo + hi) >> 1;
    if (!has_[i]) { tree_[i] = g; has_[i] = true; return; }   // empty: take it
    Line f = tree_[i];
    bool left  = g.get(lo) < f.get(lo);
    bool middle= g.get(mid) < f.get(mid);
    if (middle) swap(tree_[i], g);            // node keeps the mid-winner
    if (lo == hi) return;
    if (left != middle)                       // loser can only win on the LEFT
        insert(i<<1, lo, mid, g);
    else                                      // ... or only on the RIGHT
        insert(i<<1|1, mid+1, hi, g);
}
long long query(int i, int lo, int hi, int x){
    if (!has_[i]) return INF;
    long long res = tree_[i].get(x);
    if (lo == hi) return res;
    int mid = (lo + hi) >> 1;
    if (x <= mid) return min(res, query(i<<1, lo, mid, x));
    else          return min(res, query(i<<1|1, mid+1, hi, x));
}
\end{lstlisting}

\begin{tracebox}[Inserting three lines over the domain x=1 to 4, then reading points]
Lines: \(f_1(x)=2x+1\), \(f_2(x)=-x+6\), \(f_3(x)=x-2\), all inserted into a
\emph{minimum} container over \(x\in[1,4]\).\\
\textbf{Insert \(f_1\).} Root \([1,4]\) empty \(\Rightarrow\) root\(=f_1\).\\
\textbf{Insert \(f_2\).} Mid \(M=2\): \(f_1(2)=5\), \(f_2(2)=4\) -- \(f_2\) wins
the mid, so root\(=f_2\) and we recurse \(f_1\). At \(L=1\): \(f_1(1)=3<f_2(1)=5\)
(left better); at \(M=2\) the loser \(f_1\) is worse. left\(\ne\)mid-result
\(\Rightarrow\) \(f_1\) can only win on the \emph{left}: recurse into \([1,2]\),
which is empty, so left child\(=f_1\).\\
\textbf{Insert \(f_3\).} Root \(f_2\), \(M=2\): \(f_2(2)=4\), \(f_3(2)=0\) --
\(f_3\) wins the mid, root\(=f_3\), recurse \(f_2\). \(L=1\): \(f_2(1)=5\),
\(f_3(1)=-1\) -- \(f_2\) worse at left; mid \(f_2\) worse. \(R=4\): \(f_2(4)=2\),
\(f_3(4)=2\) tie. \(f_2\) is not better on either side \(\Rightarrow\) drop
\(f_2\) at the root's descent (it still lives in the left child from the prior
insert, which is where it mattered).\\
\textbf{Queries.} \(x{=}1\): path root\(f_3\)=-1, left child\(f_1\)=3 \(\Rightarrow\)
\(\min=-1\).  \(x{=}3\): root\(f_3\)=1, right child empty \(\Rightarrow\) \(1\).
\(x{=}4\): \(f_3(4)=2\) \(\Rightarrow\) \(2\).  The lower envelope is
\((-1,\,0,\,1,\,2)\) at \(x=1..4\), and each read is one path walk.
\end{tracebox}

\begin{observation}
The container is \emph{not} a convex hull data structure even though it maintains
an envelope: it never stores the hull vertices, only ``the line that wins at this
node's midpoint.'' That is what gives it the \(\bigO{\log C}\) worst case with
zero slope ordering, and it is why deletions are awkward (the \(\secref{c29-path}\)
workaround is a DFS \emph{rollback}, not a true delete).
\end{observation}

\begin{complexitybox}[The container]
\textbf{Insert} \(\bigO{\log C}\), \textbf{query} \(\bigO{\log C}\), each a
single root-to-leaf walk (\(\thmref{c29-lichao-cost}\)). \textbf{Space} \(\bigO{C}\)
for a pre-built tree over \([1,C]\), or \(\bigO{(\#\text{inserts})\log C}\) built
lazily. A DP that does one insert and one query per state is \(\bigO{n\log C}\)
total -- the payoff used in the problem lab.
\end{complexitybox}

% ---------------------------------------------------------------
\subsection{Problem Lab: CF 631E -- move one element, maximise the product sum}
% ---------------------------------------------------------------

\solvedproblem{c29-631e}{Product Sum}{CF 631E}
  {https://codeforces.com/problemset/problem/631/E}{CF 2600 L5}
  {Li Chao line container; DP-with-lines transition; the gain of moving element p to q rewritten as a line in the destination}

\begin{problemstatement}[CF 631E Product Sum]
Given an array \(a\) of length \(n\). Its \emph{characteristic} is
\(\sum_{i=1}^{n} i\cdot a_i\). You may perform \emph{exactly one} operation: pick
one element and move it to \emph{any} position (beginning, end, or between any two;
leaving it in place is allowed). Maximise the resulting characteristic.

\textbf{Input.} Line~1: \(n\) (\(1\le n\le 5\cdot10^{5}\)).  Line~2: the
\(n\) integers \(a_1,\dots,a_n\) (\(1\le a_i\le 10^{6}\)).

\textbf{Output.} One integer -- the maximum characteristic.

\textbf{Examples (official).}
\begin{verbatim}
Input:              Output:
4                   39
4 3 2 5

Input:              Output:
5                   49
1 1 2 7 1

Input:              Output:
3                   9
1 1 2
\end{verbatim}
\end{problemstatement}

\begin{attemptgate}[Attempt it yourself first -- \lv{5}, give yourself 40 minutes]
Write the gain of ``remove index \(p\), re-insert at index \(q\)'' as an explicit
expression in the prefix sums; then notice one of its two halves is a
line in \(q\).
\peek{how do the coefficients of the elements between \(p\) and \(q\) change when you lift \(a_p\) out and drop it at \(q\)?}
\nudge{if \(q>p\): the elements \(p{+}1..q\) shift one slot left (lose one times
their value), and \(a_p\) moves from weight \(p\) to weight \(q\).  gain\(=(q-p)a_p-(P_q-P_p)\).}
\fullhint{For \(q>p\): \(\text{gain}=(q a_p-P_q)-(p a_p-P_p)\).  Fixing the
\emph{destination} \(q\) and maximising over \(p<q\), the part that varies with
\(p\) is \(a_p\cdot q+(P_p-p\,a_p)\) -- a \emph{line in \(q\)} with slope \(a_p\)
and intercept \(P_p-p\,a_p\).  Sweep \(q=1..n\), keeping all source-lines \(p<q\)
in a Li Chao container (\(\coderef{c29-lichao}\)); at \(q\) read
\(\max_p\) and subtract \(P_q\).  The mirrored case \(q<p\) gives lines with
intercept \(P_{p-1}-p\,a_p\) queried the same way in a reverse sweep.  Best gain
over both, plus the original \(S\), is the answer -- sample~1: \(S=36\), best gain
\(3\) \(\Rightarrow\) \(39\). \(\bigO{n\log C}\).}
\end{attemptgate}

\solstep{Understanding}
The characteristic is a weighted sum with weights \(1,2,\dots,n\).  Moving one
element re-weights it and shifts the block it jumps over.  Enumerating the
\(O(n^2)\) source/destination pairs is the naive route; the Li Chao angle is to
treat the \emph{destination} as the query point \(x\) and every feasible
\emph{source} as a line, so each destination is resolved with one container query.

\solstep{Key Observations}
\begin{observation}
Moving \(a_p\) to a later \(q>p\) changes the sum by
\(\text{gain}(p,q)=(q-p)\,a_p-(P_q-P_p)\), where \(P\) is the prefix-sum array.
Rearranged in \(q\): \(\text{gain}=\big(a_p q+(P_p-p\,a_p)\big)-P_q\).
\end{observation}
\begin{observation}
For a fixed destination \(q\), the bracket is a \emph{line evaluated at \(x=q\)}:
slope \(m=a_p\), intercept \(b=P_p-p\,a_p\).  Slopes \(a_p\) are arbitrary
positive integers, so a monotone hull deque does not apply -- the Li Chao
container does (\(\defref{c29-lichao}\)).
\end{observation}
\begin{observation}
The forward (\(q>p\)) and backward (\(q<p\)) moves decouple: run one sweep for
each.  The backward gain is \(\text{gain}(p,q)=(q-p)a_p+(P_{p-1}-P_{q-1})
=\big(a_p q+(P_{p-1}-p\,a_p)\big)-P_{q-1}\), again a line queried at \(x=q\).
\end{observation}

\solstep{The Approach Ladder}
\begin{center}\renewcommand{\arraystretch}{1.25}\setlength{\tabcolsep}{3pt}
\begin{tabular}{@{}c p{0.26\linewidth} p{0.30\linewidth} l c@{}}
\toprule
\textbf{\#} & \textbf{Approach} & \textbf{Idea} & \textbf{Time} & \textbf{Verdict}\\
\midrule
1 & \lv{2}\ brute force & try every \((p,q)\), recompute sum & \(\bigO{n^3}\) & \tle \\
2 & \lv{3}\ prefix-scan pairs & every \((p,q)\) via prefix sums & \(\bigO{n^2}\) & \tle \\
3 & \lv{4}\ monotone hull & assume sorted slopes & -- & \textcolor{warningred}{\ding{55}\ \textbf{WA}} \\
4 & \lv{5}\ \textbf{Li Chao over destinations} & two sweeps, source-as-line (\(\coderef{c29-631e}\)) & \(\bigO{n\log C}\) & \ac \\
\bottomrule
\end{tabular}\end{center}
Rung~3 breaks because the slopes \(a_p\) arrive unsorted, so the hull is not
maintained -- the classic wrong-tool trap the recognition card flags.

\solstep{Solution}
Compute \(S=\sum i\,a_i\) and the prefix sums \(P\).  Sweep \(q=1..n\): at each
\(q\), if any line is present, query the container at \(x=q\) for \(V\), and
update the answer with \(V-P_q\); then insert the line for source \(p=q\)
(slope \(a_q\), intercept \(P_q-q\,a_q\)).  Do the mirror sweep \(q=n..1\)
with intercept \(P_{q-1}-q\,a_q\) and subtract \(P_{q-1}\).  Take the maximum
gain over both sweeps and print \(S+\text{best}\).

\solstep{Implementation}
\begin{lstlisting}[caption={CF 631E: forward sweep and backward sweep, each a Li Chao container of source-lines queried at the destination x=q; answer = original S + best gain}, label={lst:c29-631e}]
#include <bits/stdc++.h>
using namespace std;
typedef long long ll;
const ll NEG = -(1LL<<62);
struct L { ll m,b; ll get(ll x) const { return m*x+b; } };
int n; ll a[500005], P[500005];
L tr[4*500005]; bool has[4*500005];
void ins(int i,int lo,int hi,L g){           // MAX container: keep mid-best
    int mid=(lo+hi)>>1;
    if(!has[i]){ tr[i]=g; has[i]=true; return; }
    L f=tr[i];
    bool le=g.get(lo)>f.get(lo), mi=g.get(mid)>f.get(mid);
    if(mi) swap(tr[i],g);
    if(lo==hi) return;
    if(le!=mi) ins(i<<1,lo,mid,g); else ins(i<<1|1,mid+1,hi,g);
}
ll qry(int i,int lo,int hi,int x){
    if(!has[i]) return NEG;
    ll res=tr[i].get(x); if(lo==hi) return res;
    int mid=(lo+hi)>>1;
    return max(res, x<=mid ? qry(i<<1,lo,mid,x) : qry(i<<1|1,mid+1,hi,x));
}
int main(){
    scanf("%d",&n); ll S=0;
    for(int i=1;i<=n;i++){ scanf("%lld",&a[i]); P[i]=P[i-1]+a[i]; S+=(ll)i*a[i]; }
    ll best=0;                                  // moving to own place = gain 0
    memset(has,0,sizeof has);                   // forward: p<q, gain=V-P[q]
    for(int q=1;q<=n;q++){
        if(q>1) best=max(best, qry(1,1,n,q)-P[q]);
        ins(1,1,n,{a[q], P[q]-(ll)q*a[q]});
    }
    memset(has,0,sizeof has);                   // backward: p>q, gain=V-P[q-1]
    for(int q=n;q>=1;q--){
        if(q<n) best=max(best, qry(1,1,n,q)-P[q-1]);
        ins(1,1,n,{a[q], P[q-1]-(ll)q*a[q]});
    }
    printf("%lld\n", S+best);
}
\end{lstlisting}

\begin{tracebox}[Sample 1: \(a=(4,3,2,5)\), \(S=36\)]
Forward sweep, container of source-lines \(m=a_p,\ b=P_p-p\,a_p\); gains at each
\(q\) are \(\max_{p<q}(a_p q+b_p)-P_q\):
\(q{=}2\): line\(_1{:}\ 4x{+}0\to8\), \(-P_2{=}8{-}7=1\).
\(q{=}3\): \(\max(12,10){=}12\), \(-9=3\).
\(q{=}4\): \(\max(16,13,11){=}16\), \(-14=2\).
Backward sweep peaks at \(+3\) (source \(p{=}3\to\) dest \(q{=}1\)).  Best gain
\(=3\), so the answer is \(36+3=\boxed{39}\) -- matching the official output.  The
\(1\) and \(2\) gains correspond to moving \(a_1\) one or three steps right; the
winning \(3\) is moving \(a_1{=}4\) from index~1 to index~3 (or symmetrically
\(a_3{=}2\) back to index~1).
\end{tracebox}

\begin{warningbox}[Query \emph{then} insert, or you allow moving an element onto itself]
At destination \(q\) the container must hold only sources \(p<q\) (forward) or
\(p>q\) (backward).  If you insert the \(q\)-line before querying, gain\((q,q)=0\)
enters the maximum harmlessly, but mixing the two sweeps or inserting out of order
lets an element ``move'' using a source on the wrong side of the boundary and
over-counts the shift.  The listing keeps the strict order (query for \(q>1\)
before inserting \(q\)); mirror it in the reverse sweep.
\end{warningbox}

\begin{hardeningbox}
\begin{itemize}[nosep,leftmargin=1.4em]
  \item \textbf{Overflow.} \(n\le5\cdot10^{5}\), \(a_i\le10^{6}\): the
        characteristic and the prefix products reach \(\sim5\cdot10^{17}\) -- use
        \code{long long}, and the sentinel line at \(\pm\infty\) must be inside
        the \code{long long} range (the listing uses \(\pm2^{62}\)).
  \item \textbf{Domain vs positions.} The container is indexed by the
        \emph{destination position} \(q\in[1,n]\), so set the domain to \([1,n]\),
        not to the value range.
  \item \textbf{Both directions matter.} A best move can be backward; a forward-only
        sweep fails sample-style cases.  Run both.
  \item \textbf{``Exactly once'' } means gain~0 is always achievable (move an element
        back to its own slot), so seed \(\mathit{best}=0\).
\end{itemize}
\end{hardeningbox}

\begin{recallbox}[23.1 in one breath]
\emph{A Li Chao container is a segment tree over the \(x\)-domain; insert keeps
the mid-winning line at each node and pushes the loser to the one half where it
can still win, so insert and query are \(\bigO{\log C}\) (\(\thmref{c29-lichao-cost}\)).
Any transition \(f=\min_j m_j x+b_j\) with unsorted slopes becomes ``insert a
line per source, query at the destination'' -- as in CF~631E.}
\selftest{\item Why does the loser of a midpoint comparison recurse to \emph{one}
child and not two? \emph{(Two lines cross at most once, so past the mid-winner it
can beat it on at most one side (\figref{c29-lichao}).)}
\item In CF~631E what is the slope and intercept of a source line, and what is
the query variable? \emph{Slope \(a_p\), intercept \(P_p-p\,a_p\) (forward); query
at the destination index \(q\) (\(\coderef{c29-631e}\)).)}}
\end{recallbox}
